\input{./._relpath-to-latexroot.ltx}
\documentclass[\latexroot/\projectname]{subfiles}
\input{\latexroot/@resources/subfile-setup.ltx}

\begin{document}

\begin{table}[h]
  \centering
  \caption{Responses of labour income and inflation to government spending (2008-12)\whenintegrated{\label{tab:qa-labour-income-inflation}}}
  \begin{tabular}{lcccc}
    \hline
     & \multicolumn{2}{c}{Labour income} & \multicolumn{2}{c}{Inflation} \\
     & \multicolumn{2}{c}{QCEW} & \multicolumn{2}{c}{BLS} \\
     \cmidrule(l){2-3} \cmidrule(l){4-5}
    \textbf{Variable data source} & \textbf{OLS} & \textbf{IV} & \textbf{OLS} & \textbf{IV} \\
    \hline
    Government spending & 1.33***  & 0.81*  & -0.000  & -0.0002  \\
     & (0.34)  & (0.41) & (0.001) & (0.002) \\
    Partial $F$ stat. & - & 78.6 & - & 166.5 \\
    County controls/state fixed effects & Yes & Yes & Yes & Yes \\
    \# of counties & 2,916 & 2,916 & 1,116 & 1,116 \\
    \hline
  \end{tabular}
  \begin{flushleft}
    \footnotesize
    \textit{Notes:} The first two columns show estimates of a regression of percentage change in labour income on cumulative government spending at the county level during the period 2008-12. The last two columns show estimates of a regression of percentage point change in inflation on log-cumulative government spending at the county level during the period 2008-12. We show results for our OLS and IV specification and the standard errors in parentheses. One, two, and three stars denote significance at the $10 \%, 5 \%$, and $1 \%$ levels, respectively.
  \end{flushleft}
\end{table}

\smartbib

\end{document}
